\section{Benchmarks and Metrics}

\subsection{Benchmark Datasets}

The evaluation in this thesis uses two benchmarks. CounterFact~\cite{meng2022locating} is the primary benchmark for direct knowledge editing behavior. Each record provides an edit prompt, a set of rephrased prompts for the same fact, and locality prompts that should preserve their original answers after editing. This structure makes CounterFact suitable for measuring whether an edit is successful on the target fact while remaining sufficiently localized.

MQuAKE-Remastered~\cite{zhong2024mquake} is used to evaluate portability. Its questions are constructed so that the correct answer depends on the edited fact through downstream compositional structure. As a result, it does not only test whether the revised fact can be reproduced directly, but also whether the new fact remains usable inside later reasoning.

For the present study, both benchmarks were prepared in two prompt styles: sentence completion (cloze) and question answering (QnA). The two styles probe the same underlying edits through different surface forms and therefore do not necessarily produce identical numerical results. In the current report, the main quantitative discussion emphasizes the cloze track because it provides the cleanest common comparison across all evaluated baselines and milestones, while the QnA track is treated as a complementary evaluation view rather than the primary reporting axis.

\subsection{Evaluation Metrics}

The present study follows the practical metric grouping used in the evaluation benchmark. Although reliability and generalization remain conceptually distinct, they are reported jointly under one headline metric called \textit{success}. Accordingly, the evaluation uses three reported metrics: success, locality, and portability. This choice is not merely terminological. It reflects the fact that, in the actual benchmark, a useful edit is expected to answer both the direct prompt and its rephrases correctly, and these two behaviors are treated together as one primary dimension of performance.

Let $e$ be an edit, let $x_e$ denote its direct edit prompt, let $\mathcal{P}_e$ denote the set of rephrased prompts associated with the same fact, let $\mathcal{N}_e$ denote the locality prompt set, and let $\mathcal{R}_e$ denote the portability prompt set. For an edited model $f_{\phi}^{(t)}$, the per-edit metrics are defined as follows.

\paragraph{Success.} Success measures whether the edited fact is expressed correctly both on the direct prompt and on its paraphrased variants. Rather than separating reliability and generalization into two different reported scores, success treats them as one in-scope performance quantity:
\begin{equation}
\mathrm{Suc}(f_{\phi}^{(t)}; e)
= \frac{1}{1 + |\mathcal{P}_e|}
\left[
\mathds{1}\left\{ f_{\phi}^{(t)}(x_e) \in \mathcal{Y}_e^* \right\}
+
\sum_{x \in \mathcal{P}_e}
\mathds{1}\left\{ f_{\phi}^{(t)}(x) \in \mathcal{Y}_e^* \right\}
\right].
\end{equation}
This definition makes the intended meaning explicit. The direct edit prompt and the rephrased prompts are not two unrelated criteria; they are two views of the same edited fact. A method that succeeds only on the original prompt but fails on paraphrases has not yet established a robust factual update, and the success metric reflects that by averaging across the full in-scope prompt set.

\paragraph{Locality.} Locality measures whether prompts outside the edited fact preserve their original answers. Let $f_{\phi}$ denote the original model before editing. The per-edit locality score is defined by
\begin{equation}
\mathrm{Loc}(f_{\phi}^{(t)}, f_{\phi}; e)
= \frac{1}{|\mathcal{N}_e|}
\sum_{x \in \mathcal{N}_e}
\mathds{1}\left\{ f_{\phi}^{(t)}(x) = f_{\phi}(x) \right\}.
\end{equation}
This metric is essential because strong success alone is not sufficient. A method may enforce the revised fact aggressively while also damaging unrelated neighboring behavior. Locality is therefore treated as a direct measure of collateral interference instead of a secondary diagnostic quantity.

\paragraph{Portability.} Portability measures whether the revised fact can be used in downstream questions whose answers depend on it. For the portability set $\mathcal{R}_e$ associated with edit $e$, the per-edit score is
\begin{equation}
\mathrm{Por}(f_{\phi}^{(t)}; e)
= \frac{1}{|\mathcal{R}_e|}
\sum_{(x^r, y^r) \in \mathcal{R}_e}
\mathds{1}\left\{ f_{\phi}^{(t)}(x^r) = y^r \right\}.
\end{equation}
Portability is the most demanding of the three metrics because it evaluates not only whether a fact can be retrieved, but whether that fact can participate in later inference. For this reason, portability values are often lower than direct success values even when the edit itself appears locally correct.

\subsection{Lifelong Aggregation Across the Edit Stream}

The previous definitions describe the behavior of one edit at one time step. Lifelong model editing, however, evaluates cumulative behavior after repeated updates. Let $\mathcal{E}_t = \{e_1, e_2, \ldots, e_t\}$ denote the set of edits that have been applied up to step $t$. The aggregate success at step $t$ is defined as
\begin{equation}
\mathrm{Suc}(f_{\phi}^{(t)}; \mathcal{E}_t)
= \frac{1}{t}
\sum_{i=1}^{t}
\mathrm{Suc}(f_{\phi}^{(t)}; e_i).
\end{equation}
This quantity captures the central difficulty of lifelong editing. Even if the newest edit is successful, the aggregate score can still decline if earlier edits have been forgotten or distorted by later updates.

Aggregate locality is defined analogously:
\begin{equation}
\mathrm{Loc}(f_{\phi}^{(t)}, f_{\phi}; \mathcal{E}_t)
= \frac{1}{t}
\sum_{i=1}^{t}
\mathrm{Loc}(f_{\phi}^{(t)}, f_{\phi}; e_i).
\end{equation}
Although locality is evaluated through per-edit neighborhoods, the aggregate value measures whether the edited model continues to preserve unrelated behavior across the full cumulative edit history.

Aggregate portability is defined by
\begin{equation}
\mathrm{Por}(f_{\phi}^{(t)}; \mathcal{E}_t)
= \frac{1}{t}
\sum_{i=1}^{t}
\mathrm{Por}(f_{\phi}^{(t)}; e_i).
\end{equation}
This aggregate is particularly informative in the lifelong setting because it reveals whether newly inserted knowledge remains usable after repeated interference, not only whether a single edit can temporarily support downstream reasoning.

Taken together, these aggregate definitions align the metric layer with the actual lifelong protocol used in the experiments. The edited model is evaluated not only on the newest incoming fact, but on the cumulative state that remains after many successive edits.
